\documentclass[11pt,a4paper]{article}
\usepackage[T1]{fontenc}
\usepackage[utf8]{inputenc}
\usepackage[main=italian,provide=*]{babel}
\usepackage{amsmath,amssymb}
\usepackage{geometry}
\geometry{margin=2.3cm,top=2.5cm,bottom=2.7cm}
\usepackage{booktabs}
\usepackage[table]{xcolor}
\usepackage{tcolorbox}
\tcbuselibrary{skins,breakable}
\usepackage{titlesec}
\usepackage{enumitem}
\usepackage{seqsplit}
\usepackage{needspace}
\usepackage{fancyhdr}
\usepackage{hyperref}
\hypersetup{colorlinks=true,linkcolor=Sepia,citecolor=Sepia,urlcolor=Sepia}

\definecolor{inkblue}{HTML}{1B3A5C}
\definecolor{lightblue}{HTML}{EAF1F8}
\definecolor{accent}{HTML}{B0562A}
\definecolor{softgray}{HTML}{6B6B6B}
\definecolor{okgreen}{HTML}{2E6B3E}
\definecolor{okgreenbg}{HTML}{EAF5EC}
\definecolor{warnamber}{HTML}{9C6B12}
\definecolor{warnbg}{HTML}{FBF1E0}
\definecolor{advisorpurple}{HTML}{5B3A7A}
\definecolor{advisorbg}{HTML}{F1ECF6}
\definecolor{notebar}{HTML}{9C6B12}

\titleformat{\section}
  {\normalfont\Large\bfseries\color{inkblue}}
  {\thesection}{0.6em}{}[{\color{inkblue!40}\titlerule}]
\titlespacing{\section}{0pt}{0.8em}{0.4em}
\titleformat{\subsection}
  {\normalfont\large\bfseries\color{accent}}
  {}{0em}{}
\titlespacing{\subsection}{0pt}{0.5em}{0.2em}

\pagestyle{fancy}
\fancyhf{}
\renewcommand{\headrulewidth}{0.4pt}
\fancyhead[R]{\small\color{softgray}\thepage}
\fancyfoot[C]{\small\color{softgray}Tesi triennale, Samuele Schiavi --- Relatore: Prof.\ Alessandro Chiesa}

\newtcolorbox{keyeq}[1][]{
  colback=lightblue, colframe=inkblue!70, boxrule=0.6pt,
  arc=2pt, left=8pt, right=8pt, top=4pt, bottom=4pt, #1
}
\newtcolorbox{panelbox}[1]{
  colback=white, colframe=softgray!50, boxrule=0.5pt,
  arc=2pt, left=7pt, right=7pt, top=3pt, bottom=3pt,
  fonttitle=\bfseries\color{inkblue}, title={#1}
}
\newtcolorbox{okbox}[1]{
  colback=okgreenbg, colframe=okgreen!60, boxrule=0.6pt,
  arc=2pt, left=7pt, right=7pt, top=3pt, bottom=3pt,
  fonttitle=\bfseries\color{okgreen}, title={#1}
}
\newtcolorbox{warnbox}[1]{
  colback=warnbg, colframe=warnamber!70, boxrule=0.6pt,
  arc=2pt, left=7pt, right=7pt, top=3pt, bottom=3pt,
  fonttitle=\bfseries\color{warnamber}, title={#1}
}
\newtcolorbox{advisorbox}[1]{
  colback=advisorbg, colframe=advisorpurple!60, boxrule=0.6pt,
  arc=2pt, left=7pt, right=7pt, top=4pt, bottom=4pt,
  fonttitle=\bfseries\color{advisorpurple}, title={#1}
}
\newtcolorbox{editnote}{
  colback=white, colframe=white, boxrule=0pt,
  borderline west={2.2pt}{0pt}{notebar},
  arc=0pt, left=10pt, right=4pt, top=2pt, bottom=2pt,
  fontupper=\itshape\small\color{softgray!90!black}
}
\fancyhead[L]{\small\color{softgray}Note su \texttt{quantum\_simulation\_dimero\_trotter.ipynb}}

\title{\vspace{-1.5em}\color{inkblue}Note di approfondimento\\[0.1em]\texttt{quantum\_simulation\_dimero\_trotter.ipynb}}
\author{Note di lavoro --- Tesi triennale, Samuele Schiavi\\ Relatore: Prof.\ Alessandro Chiesa}
\date{}

\begin{document}
\sloppy
\maketitle
\thispagestyle{fancy}
\vspace{-2em}

Materiale a corredo del notebook, tenuto fuori dalle celle per
mantenerlo snello.

\section{Cosa vuol dire "oscillazioni non monocromatiche"}

Il termine viene dall'ottica: una luce monocromatica ha una sola
frequenza (un solo colore). Un'oscillazione \textbf{monocromatica} è un
segnale con una sola frequenza pura:
$$f(t)=A\cos(\omega t+\phi)$$
Un'oscillazione \textbf{non monocromatica} è invece una sovrapposizione
di più frequenze diverse:
$$f(t)=A_1\cos(\omega_1 t+\phi_1)+A_2\cos(\omega_2 t+\phi_2)+\dots$$

\subsection*{Applicazione al dimero}

$\langle S_z\rangle(t)$ è per costruzione una somma di modi a frequenze
di Bohr $\Delta_{kl}=E_l-E_k$ (sezione 3 del notebook):
\begin{keyeq}
\centering
$\langle S_z\rangle(t)=\text{costante}+\displaystyle\sum_{k<l}2\,\mathrm{Re}\!\left[c_k^*c_l(S_z)_{kl}\,e^{-i\Delta_{kl}t}\right]$
\end{keyeq}

\begin{itemize}[leftmargin=1.3em,itemsep=0.2em]
\item Se \textbf{un solo} termine domina, il grafico è un coseno pulito
e regolare. È il caso del regime R0 ($D/J=1/6$, rapporto
$a_2/a_1=0.04$).
\item Se \textbf{due o più} termini hanno ampiezze confrontabili e
frequenze diverse, si sommano e producono \textbf{battimenti}: massimi
e minimi irregolari, senza periodicità semplice. È il caso di R1
($a_2/a_1=0.74$).
\end{itemize}

\subsection*{Perché il relatore lo chiede}

Un'oscillazione monocromatica equivale dinamicamente a un sistema a due
livelli, e non mette alla prova né il codice Trotter né lo sviluppo su
più autostati --- un singolo coseno si ottiene anche a mano. Un regime
non monocromatico costringe la dinamica su più autostati a entrare in
gioco.

\begin{editnote}
\textbf{Precisazione da tenere presente} (sez.~3.1 del notebook):
poiché $|T_0\rangle$ è esattamente disaccoppiato, la dinamica è
effettivamente a \textbf{3 livelli}, con solo \textbf{2 frequenze
indipendenti} (la terza è la somma delle altre due, vincolo automatico
fra gap di livelli). Il passaggio R0 $\to$ R1 è quindi da 1 a 2
frequenze indipendenti, non da 2 a 4 livelli: da riportare con questa
precisione, perché il guadagno è reale ma più modesto di quanto
suggerirebbe il conteggio dei livelli.
\end{editnote}

\section{Criterio di selezione del regime: cosa è stato usato e cosa no}

Nel notebook la scelta del punto di lavoro usa due grandezze
\textbf{trasparenti e verificabili a mano}, entrambe calcolate dalla
dinamica esatta:

\begin{enumerate}[leftmargin=1.4em,itemsep=0.2em]
\item \textbf{escursione} picco-picco di $\langle S_z\rangle(t)$ --- il
segnale deve essere leggibile sopra il rumore statistico
$\sim1/\sqrt{N_\text{shots}}$;
\item \textbf{$a_2/a_1$} --- rapporto fra le ampiezze dei due modi
dominanti, $\approx0$ per un segnale monocromatico, $\approx1$ per due
frequenze equilibrate.
\end{enumerate}

Servono entrambe: un segnale grande ma monocromatico non soddisfa la
richiesta del relatore (caso della risonanza $b=-J$), e un segnale
spettralmente ricco ma minuscolo è inutilizzabile.

\begin{editnote}
\textbf{Nota storica.} In una versione precedente il criterio era
un'\emph{entropia spettrale pesata} $A\cdot S$ con
$S=-\sum_m p_m\ln p_m$. È stata rimossa: il concetto generale esiste in
signal processing (Kapur \& Kesavan 1992, entropia di Shannon su una
distribuzione di potenza spettrale), ma la sua applicazione a un
insieme discreto di ampiezze di transizione fra autostati non risulta
da alcuna fonte, né dal materiale del corso --- era una proposta
originale non richiesta. Inoltre risultava poco robusta: lo scarto di
punteggio fra 1° e 5° classificato era del $4\%$, e un criterio
alternativo rimescolava completamente la graduatoria. Il criterio a due
colonne attuale è più semplice, altrettanto efficace e direttamente
verificabile dal relatore.
\end{editnote}

\needspace{10\baselineskip}
\section{Nota sulla convenzione di segno delle frequenze}

La convenzione standard in letteratura (Cohen-Tannoudji e affini) è
$$\omega_{kl}\equiv\frac{E_k-E_l}{\hbar}$$
cioè il primo indice porta il segno positivo. Nel notebook si usa
invece
\begin{keyeq}
\centering
$\Delta_{kl}\equiv E_l-E_k>0 \qquad (k<l)$
\end{keyeq}
con indici esplicitamente ordinati, per avere frequenze positive senza
dover ricordare la convenzione di segno a ogni passaggio:
$\Delta_{kl}=\omega_{lk}$. La scelta non incide sul risultato fisico
--- $\langle S_z\rangle(t)$ è reale e ogni modo entra come
$2|z|\cos(\Delta_{kl}t-\arg z)$, quindi invertire il segno di una
frequenza si riassorbe nella fase --- ma va tenuta coerente nel resto
della tesi.

\section{Convenzioni da dichiarare insieme}

Le tre convenzioni che, prese singolarmente, sembrano innocue e che
insieme producono la maggior parte degli errori di segno difficili da
diagnosticare:

\begin{enumerate}[leftmargin=1.4em,itemsep=0.3em]
\item \textbf{Endianness di Qiskit}: \texttt{Pauli('AB')} ha $B$ sul
qubit 0. Con spin 1 $\to$ qubit 0, quindi $X_1Z_2\equiv$
\texttt{Pauli('ZX')}.
\item \textbf{Segno del termine DM}: con $H_2'=-H_2$ cambiano segno
\emph{entrambi} gli angoli dei due $R_{ZZ}$, la struttura del circuito
resta identica. Per l'osservabile $S_z$ da $|00\rangle$ il risultato
non cambia (i due casi sono legati dallo scambio $1\leftrightarrow2$,
che lascia invarianti sia $|00\rangle$ sia $S_z$).
\item \textbf{Fattori $1/2$ dei gate}: $R_z(\phi)=e^{-i\phi Z/2}$,
$R_{ZZ}(\phi)=e^{-i\phi ZZ/2}$ --- da cui $\theta=Jt/2$ e non $Jt/4$
nei blocchi di scambio.
\end{enumerate}

\needspace{14\baselineskip}
\section{Provenienza del punto di lavoro R1 --- cronologia}

Vale la pena registrarla, perché il punto è cambiato due volte e il
relatore potrebbe chiedere conto del percorso.

\begin{enumerate}[leftmargin=1.4em,itemsep=0.3em]
\item \textbf{Prima versione (scartata).} R1
$=(b/J=-0.15,\;D/J=1)$ veniva da uno scan automatico su 2349 punti
classificati con un'\emph{entropia spettrale pesata} $A\cdot S$.
Criterio poco robusto (scarto 1°--5° classificato del 4\%) e non
riconducibile ad alcuna fonte: rimosso.
\item \textbf{Seconda versione (scartata).} Stessa coppia di valori, ma
giustificata con una tabella di regimi candidati. Era una
razionalizzazione \emph{a posteriori}: la tabella era stata costruita
con R1 già dentro, quindi non permetteva di ricostruire come ci si
fosse arrivati.
\item \textbf{Versione attuale.} R1 $=(b/J=-0.18,\;D/J=1)$
\textbf{derivato} dalla struttura a V dell'Hamiltoniana (sez.~3.1--3.4
del notebook): due condizioni indipendenti (ampiezza del segnale dalla
formula di Rabi; numero di frequenze dai due detuning
$\Delta_\pm=J\pm b$) e due scan 1D. I due modi dominanti risultano
bilanciati a $a_2/a_1=0.995$.
\end{enumerate}

Il valore $-0.15$ resta comunque dentro il plateau
$b/J\in[-0.22,-0.14]$: se serve mantenerlo per coerenza con materiale
già scritto, basta cambiarlo nel setup senza invalidare la derivazione.

\section{Bibliografia: cosa è citabile e cosa no}

La sezione 8 del notebook contiene la tabella completa. In sintesi, per
la derivazione di R1 \textbf{non è stata consultata alcuna fonte
esterna}: la struttura è stata ricavata e verificata numericamente. Gli
ingredienti sono però standard, e due sono già documentati nel progetto
stesso (\texttt{simmetrie\_dmi\_spin.pdf} \S 11.1, \S 11.2.2, \S 12,
\S 13.3).

\begin{editnote}
L'unico elemento senza riferimento preciso è la formula di Rabi per il
trasferimento di popolazione a due livelli: è un risultato da manuale,
ma non è stata verificata una citazione specifica e \textbf{non va
inventata}. Se serve in tesi, va presa da un testo che hai
effettivamente in mano.
\end{editnote}

La combinazione dei due criteri in una procedura di selezione è
costruzione originale di questo lavoro, e va presentata come tale.

\needspace{8\baselineskip}
\section{Discrepanza aperta con \texttt{simmetrie\_dmi\_spin.pdf}}

\begin{warnbox}{Da chiarire con il relatore}
Vedi sez.~8.1 del notebook per la verifica numerica completa. In
breve: gli elementi di matrice della DMI riportati nel documento
(\S 11.2.1 e \S 11.2.2) risultano \textbf{doppi} rispetto al calcolo
diretto. Il controllo di covarianza rotazionale
($\|(\vec s_1\times\vec s_2)_a|0,0\rangle\|=1/2$ per ogni asse $a$)
supporta il valore ricavato qui. Se confermato, va corretta anche
l'espressione di \S 13.2 per $E_\pm$. Da chiarire con il relatore
prima di portare l'una o l'altra versione in tesi.
\end{warnbox}

\end{document}
